\section{Repaso Quiz}

\subsection{Preguntas de verdadero o falso}

\begin{enumerate}
	\item Un agente que percibe solo información parcial acerca del estado no puede ser perfectamente racional: \textbf{Falso} la racionalidad depende de la expectativa de rendimiento dada la información que el agente posee (el historial de percepciones), no del éxito omnisciente o del estado real completo del mundo. Un agente en un entorno parcialmente observable puede tomar la decisión óptima posible con base en sus percepciones parciales y seguir siendo perfectamente racional.
	\item Existen ambientes de tareas en lso que ningun agente de reflejo puro se puede comportar de forma racional: \textbf{Verdadero} en entornos parcialemnte observables, las acciones optimas suelen depender de la historia o de estados ocultos que no se peuden deducir unicamente de la percepcion actual; un agente de reflejo puro carece de memoria y puede entrar en bucles infinitos.
	\item Es posible para un agente dado ser perfectamente racional en dos ambientes de tarea diferentes: \textbf{Verdadero}, si las acciones requeridas para maximizar la medida de rendimiento en ambos entornos coinciden ante las mismas percepciones, el agente será racional en ambos.
	\item Todos los agentes son racionales en un ambiente no observable: \textbf{Falso}, que el entorno sea no observable no significa que cualquier acción sea buena. El agente debe actuar basandose en estimaciones o en su conocimiento inicial (suponer estados). Un agente que tome decisiones autodestructivas a ciegas seguirá siendo irracional, mientras que uno racional calculará la secuencia de acciones que estadisticamente tenga mayor probabilidad de exito.
	\item Un agente de poker perfectamente racional nunca pierde: \textbf{Falso}, el poker es un entorno estocastico y de informacion oculta, un agente perfectamente racional tomara la decisión con \textit{el mayor valor esperado}, pero la suerte puede hacer que pierda manos o partidas.
	
\end{enumerate}

\subsection{Descripción PEAS y caracterizacion del ambiente:  Rutina de piso de gimnasia}

\textbf{P (Rendimiento/performance)}: Puntuación de los jueces (dificultad, ejecución, fluidez, sincronización, técnica, no salirse del tapete)
\textbf{E (Entorno/Environment)}: El tapete de gimnasia, la gravedad, las capacidades fisicas del cuerpo del gimnasta.
\textbf{A (Actuadores/Actuatprs)}: Musculos (piernas, brazos, torso) para realizar giros, saltos, mortales y transiciones.
\textbf{S (Sensores/Sensors)}: Sistema vestibular (equilibrio), propiocepción, vision.

\subsection{Caracterización del ambiente: Rutina de piso de gimnasia}

\textit{Continuo o discreto?} \textbf{Continuo} ya que involucra las posiciones corporales, fuerzas y tiempos que varían continuamente.

\textit{Total o parcialmente observable?} \textbf{Totalmente observable} ya que el agente conoce su propio estado fisico y la posición del piso.

\textit{Determinista o estocástico?} \textbf{Estocástico} ya que a nivel fisico real, pequeñas variaciones o fatiga muscular pueden alterar el resultado de un salto.

\textit{Episodico o Secuencial?} \textbf{Secuencial} un mal salto inicial arruina el equilibrio para el siguiente movimiento.

\textit{Estatico o dinámico?} \textbf{Estático} ya que el tapete no cambia de lugar por si mismo mientras se ejecuta la rutina, aunque el tiempo corre.

\textit{Un agente o multiagente?} \textbf{agente} ya que es una rutina individual.


\subsection{Descripción PEAS y caracterizacion del ambiente:  Jugar Futbol}

\textbf{P (Rendimiento/performance)}: Goles anotados, evitar goles en contra, ganar el partido, precisión de pases, posesión del balón.
\textbf{E (Entorno/Environment)}: Campo de juego, 11 rivales, 10 compañeros de equipo, condiciones climaticas.
\textbf{A (Actuadores/Actuatprs)}: Movimiento de las piernas(correr, patear, cambiar de dirección), torso, cabeza.
\textbf{S (Sensores/Sensors)}: Visión, audición, tacto/propiocepción.

\subsection{Caracterización del ambiente: Jugar futbol}

\textit{Continuo o discreto?} \textbf{Continuo} ya que involucra las posiciones corporales, fuerzas y tiempos que varían continuamente.

\textit{Total o parcialmente observable?} \textbf{parcialmente observable} ya que no se puede ver lo que pasa a la espalda del jugador sin voltear.

\textit{Determinista o estocástico?} \textbf{Estocástico} fisica del balo, rebotes impredecibles.

\textit{Episodico o Secuencial?} \textbf{Secuencial} 

\textit{Estatico o dinámico?} \textbf{Dinámico} ya que los rivales y compañeros se mueven continuamente mientras el agente decide.

\textit{Un agente o multiagente?} \textbf{multiagente} ya que es competitivo con los rivales y cooperativo con los compañeros.


\section{Formulación de problemas de busqueda}

\begin{enumerate}
	\item \texttt{Por qué la formulación debe estar guiada por una formulacion de la meta?} En un entorno real, los detalles del mundo son infinitos. Si no se define primero un \textbf{objetivo} claro, el agente no sabrá que aspectos del mundo son relevantes y cuales debe ignorar \textit{(abstraccion)}. la meta sirve como criterio para establecer que constituye un estado satisfactorio, lo que permite acotar las acciones posibles a considerar y estructurar la función de costo y la prueba de meta. Sin ella el espacio de estados sería inmanejable.
\end{enumerate}


